\chapter{L'apparato scheletrico: generalità e cranio}

% ---------------------------------------------------------------------------
\section{Funzioni dell'apparato scheletrico}

\textbf{Apparato scheletrico}: \textbf{protezione}, \textbf{forma} e \textbf{sostegno},
\textbf{movimento}, \textbf{riserva di minerali}, \textbf{emopoiesi}. Si distingue in
\textbf{appendicolare} e \textbf{assile}.
\begin{itemize}
  \item $\sim$20\% del peso corporeo; nell'adulto $\sim$206 ossa (esclusi ossicini dell'udito e
    sesamoidi);
  \item conformazione modellata dalla \textbf{forza di gravità} e dall'uso di arti inferiori e
    superiori.
\end{itemize}

% ---------------------------------------------------------------------------
\section{Evoluzione della postura eretta}

Passaggio da \textbf{quadrupede} a \textbf{bipede} $\rightarrow$ rimodellamento del \textbf{bacino}:
\begin{itemize}
  \item mammiferi inferiori: bacino stretto sul piano frontale, allungato in senso sagittale;
  \item uomo: asse laterolaterale maggiore, asse sagittale minore.
\end{itemize}
Evoluzione del cranio: scimpanzé $\rightarrow$ \textit{Australopithecus africanus} (3,7 mln di anni)
$\rightarrow$ \textit{Homo abilis} (2,2 mln) $\rightarrow$ \textit{Homo erectus} (1 mln) $\rightarrow$
\textit{Homo sapiens} (da 200.000 anni) $\rightarrow$ aumento della capacità cranica e riduzione del
massiccio facciale.

% ---------------------------------------------------------------------------
\section{Equilibrio corporeo e arti}

\begin{itemize}
  \item \textbf{arti inferiori}: deambulazione ed equilibrio statico/dinamico; più robusti dei
    superiori;
  \item \textbf{arti superiori}: equilibrio, deambulazione e varietà di azioni;
  \item la \textbf{colonna vertebrale} presenta curvature opposte (di lato) $\rightarrow$
    riequilibrio del peso e ammortizzazione della gravità su cranio e tronco.
\end{itemize}

\rubrica{Veduta anteriore dello scheletro}
\begin{itemize}
  \item \textbf{cranio}: orbita, mascella, mandibola, articolazione temporo-mandibolare;
  \item \textbf{cingolo scapolare}: clavicola (articolazione sternoclavicolare) e scapola (processo
    coracoideo, acromion);
  \item \textbf{arto superiore}: omero (piccola e grande tuberosità); avambraccio = radio + ulna
    (articolazioni radioulnari prossimale e distale); mano = carpo, metacarpo, dita;
  \item \textbf{sterno}: manubrio, corpo, processo xifoideo;
  \item \textbf{bacino}: iliaco, pubico, ischiatico (sinfisi pubica; sacro e coccige); anca
    (coxo-femorale);
  \item \textbf{arto inferiore}: femore; ginocchio (femoro-patellare, patella); gamba = tibia
    (tuberosità tibiale) + perone; piede = astragalo, ossa tarsali, navicolare, cuneiformi, cuboide,
    metatarso, falangi.
\end{itemize}

\rubrica{Veduta posteriore}
\begin{itemize}
  \item cranio: parietale e occipitale; atlante ed epistrofeo;
  \item spalla: spina della scapola, articolazione acromioclavicolare, testa dell'omero nella
    cavità glenoidea;
  \item colonna vertebrale; gomito (olecrano dell'ulna); cresta iliaca; coxofemorale (collo del
    femore, acetabolo);
  \item carpo: scafoide, semilunare, piramidale, uncinato, capitato, trapezoide, trapezio;
  \item femore: testa, piccolo e grande trocantere; ginocchio: condilo mediale e laterale;
  \item gamba/caviglia: articolazione tibioperoneale, testa del perone, piatto tibiale, malleolo
    mediale, sindesmosi tibioperoneale, malleolo laterale, articolazione talocrurale, astragalo,
    subtalare, calcagno.
\end{itemize}

% ---------------------------------------------------------------------------
\section{Il cranio}

\begin{itemize}
  \item articolazioni tra le ossa craniche/facciali = tutte \textbf{sinartrosi} (per lo più
    \textbf{suture}); una sola mobile = \textbf{diartrosi} tra mandibola e osso temporale;
  \item \textbf{neurocranio}: protegge l'encefalo;
  \item \textbf{splancnocranio} (massiccio facciale): in rapporto con la porzione iniziale di
    apparato respiratorio e digerente.
\end{itemize}

\subsection{Neurocranio e splancnocranio}
\begin{itemize}
  \item \textbf{neurocranio}: \textbf{base} e \textbf{volta}, ossa piatte e brevi $\rightarrow$
    occipitale, sfenoide, temporale, etmoide, frontale, parietale;
  \item \textbf{splancnocranio}: lacrimale, nasale, zigomatico, mascellare, mandibola.
\end{itemize}
Suture principali: sagittale (tra i parietali), lambdoidea (parietali--occipitale), coronale
(parietali--frontale), squamosa (parietale--temporale). Elementi della base cranica (visione
inferiore, mandibola rimossa): forame palatino maggiore, fossa incisiva, palatino, vomere, processi
pterigoidei dello sfenoide, arco zigomatico, forame ovale, forame lacero, canale carotideo, meato
acustico esterno, fossa mandibolare, forame giugulare, forame stilomastoideo, forame magno, linea
nucale superiore.

\subsection{Osso occipitale}
Impari, a conchiglia/romboidale, base del neurocranio.
\begin{itemize}
  \item \textbf{forame magno} ($\phi$ medio $\sim$35 mm): passaggio di midollo spinale, arterie
    vertebrali, radici spinali, XI paio dei nervi cranici; il bulbo vi si estende nel canale
    vertebrale $\rightarrow$ midollo spinale;
  \item si articola anteriormente per sincondrosi (via clivo) col corpo dello sfenoide.
\end{itemize}

\subsection{Osso sfenoide}
Impari e mediano, al centro della base cranica.
\begin{itemize}
  \item articolazioni: etmoide e frontale (avanti), occipite (dietro), temporale e parietale (lati);
  \item partecipa a orbite, fosse nasali, fossa temporale e sotto-temporale;
  \item corpo cuboidale a 6 facce (2 laterali + superiore endocraniche; anteriore e inferiore
    esocraniche; posteriore per sinostosi col clivo); contiene i \textbf{seni sfenoidali}
    (comunicanti con le cavità nasali).
\end{itemize}

\subsection{Osso temporale}
Pari; anteriore all'occipitale, posteriore alle grandi ali dello sfenoide, inferiore al parietale.
\begin{itemize}
  \item parte interna = \textbf{rocca petrosa} (mediale e in avanti); parte esterna = \textbf{mastoidea} (in basso);
  \item presenta processo zigomatico, fossa glenoidea, processo stiloideo.
\end{itemize}

\subsection{Osso etmoide}
Impari e mediano; parte delle cavità nasali e orbitarie.
\begin{itemize}
  \item articolazioni: frontale (alto), sfenoide (dietro), mascellari e vomere (avanti/basso),
    lacrimali (lati);
  \item due lamine perpendicolari + due masse laterali;
  \item \textbf{lamina cribrosa} (orizzontale): forellini per il nervo olfattivo; faccia superiore
    = parte della fossa cranica, inferiore = volta delle cavità nasali; porta la \textbf{crista
    galli};
  \item masse laterali = labirinto etmoidale; lamina perpendicolare = parte del setto nasale.
\end{itemize}

\subsection{Osso frontale}
Impari e mediano, costituisce la volta cranica.
\begin{itemize}
  \item articolazioni: parietali, etmoide, mascellari, nasali, sfenoide, lacrimali;
  \item parte verticale (\textbf{squama}) + parte orizzontale; margini superiore (parietale),
    inferiore (sopraorbitario), posteriore (sfenoidale);
  \item elementi: sutura metopica, arcata sopracciliare, margine e foro sopraorbitario, cresta
    frontale, incisura etmoidale, seni frontali, fossa orbitaria.
\end{itemize}

\subsection{Osso parietale}
Pari, piatto, quadrangolare irregolare, concavo verso l'interno.
\begin{itemize}
  \item articolazioni: controlaterale (sutura sagittale), frontale (avanti), grande ala dello
    sfenoide e squama del temporale (lati), occipitale (dietro);
  \item faccia esocranica convessa (eminenza parietale); faccia endocranica concava con solchi
    dell'arteria meningea media; linee temporali superiore e inferiore.
\end{itemize}

% ---------------------------------------------------------------------------
\section{Lo splancnocranio}

\subsection{Osso nasale}
Piccolo, laminare, pari, trapezoidale. Articolazioni: frontale (alto), controlaterale (mediale),
processo frontale della mascella (laterale); margine inferiore = contorno superiore dell'apertura
piriforme.

\subsection{Osso lacrimale}
Laminare, pari, quadrangolare. Con la mascella forma la fossa del sacco lacrimale.

\subsection{Osso palatino}
Laminare, pari, tra processo pterigoideo dello sfenoide e superficie posteromediale del mascellare.
Partecipa a cavità nasali, palato e (in parte) pavimento dell'orbita. Porzione verticale +
orizzontale.

\subsection{Cornetto nasale inferiore}
Pari, sotto il cornetto etmoidale medio, nelle cavità nasali; lamina a tegola allungata. Tre
processi superiori:
\begin{itemize}
  \item \textbf{etmoidale}: si unisce al processo uncinato dell'etmoide;
  \item \textbf{lacrimale}: con l'osso lacrimale, chiude il canale nasolacrimale;
  \item \textbf{mascellare}: con la parte inferiore dello iato mascellare.
\end{itemize}

\subsection{Vomere}
Lamina quadrangolare impari, parte del setto nasale. Articolazioni: sfenoide (alto/dietro), lamina
perpendicolare dell'etmoide (alto/avanti), parte orizzontale del palatino e processi palatini delle
mascelle (basso).

\subsection{Osso zigomatico}
Osso breve, pari, quadrangolare.
\begin{itemize}
  \item processo frontale (angolo superiore) $\rightarrow$ osso frontale;
  \item angoli anteriore/inferiore $\rightarrow$ mascella;
  \item processo temporale (angolo posteriore) $\rightarrow$ osso temporale = arco zigomatico;
  \item bordo postero-inferiore $\rightarrow$ grande ala dello sfenoide.
\end{itemize}

\subsection{Osso mascellare}
Pari, il più grande del massiccio facciale; coi controlaterali forma l'arcata mascellare superiore.
\begin{itemize}
  \item corpo + 4 processi: zigomatico, frontale, alveolare, palatino;
  \item corpo piramidale a 4 facce: anteriore (sporgenze sulle radici dei denti, incisura nasale),
    infratemporale/posteriore (separata dall'anteriore dal processo zigomatico).
\end{itemize}

\subsection{Mandibola}
Impari e mediano.
\begin{itemize}
  \item corpo a ferro di cavallo (si articola coi denti dell'arcata inferiore, molari inclusi);
  \item due rami $\rightarrow$ osso temporale; ciascun ramo termina col processo condiloideo
    (condilo) e coronoideo, separati dall'incisura mandibolare;
  \item faccia interna: rilievi/fossette per muscoli estrinseci della lingua e sopraioidei (linea
    miloioidea); foro del canale mandibolare (nervo mandibolare), che emerge dal forame mentoniero
    (faccia esterna, presso la protuberanza mentoniera);
  \item margini: alveolare e basale.
\end{itemize}

% ---------------------------------------------------------------------------
\section{Osso ioide}

\textbf{Osso ioide}: impari e mediano, a ferro di cavallo con concavità posteriore, nella regione
anteriore del collo (in rapporto con cartilagine tiroidea, membrana tiro-ioidea, trachea, muscoli
del pavimento della bocca).
\begin{itemize}
  \item unico osso che \textbf{non si articola} con altre ossa: posizione mantenuta solo dai
    muscoli inseriti;
  \item muscoli: digastrico, stiloioideo, tireoioideo, omoioideo, sternoioideo, costrittore medio
    della faringe, ioglosso, condroglosso, genioglosso, genioioideo.
\end{itemize}
